\documentclass[11pt]{article}

\usepackage[margin=1in]{geometry}
\usepackage{amsmath}
\usepackage{amssymb}
\usepackage[round]{natbib}
\usepackage{hyperref}

% \H is already defined by the LaTeX kernel as the Hungarian umlaut accent, so it
% must be redefined rather than newly defined.
\newcommand{\X}{\mathcal{X}}
\newcommand{\Y}{\mathcal{Y}}
\renewcommand{\H}{\mathcal{H}}
\newcommand{\Hset}{\mathbb{H}}

\title{Notes on Bayesian Optimization of Composite Functions}
\date{\today}

\begin{document}

\maketitle

Bayesian optimization is a class of methods for optimizing black-box functions $y(x)$.
Bayesian optimization of composite functions proposes that optimization can be accelerated through so-called ``grey-box information''.  Specifically, it considers settings where $y(x) = g(h(x))$ for some known function $g$ and some black-box vector-valued function $h(x)$.  \citet{pmlr-v97-astudillo19a} showed that performance can be significantly accelerated by fitting a Gaussian process on $h$, which then implies a posterior distribution on $y(x)$, and then using expectetd improvement with respect to this new posterior distribution on $y(x)$.

Significant improvements were observed in settings where $g(z) = \sum_{j=1}^m z_j^2$ but unpublished experiments on applications using other functions $g(\cdot)$ have seen much smaller improvements.

This paper tackles the question, {\bf when does the composite modeling provide significant benefits?}

\section{Model}

We consider a setting where we wish to predict and possibly optimize a function $y(x)$,
We suppose that $y(x) = g(h(x))$ where $x \in \mathbb{R}^d$ is a vector, $h(x) \in \mathbb{R}^m$ is an unknown vector-valued function whose output we can observe when performing experiments, $g(z)$ is a known function that takes $z \in \mathbb{R}^m$ as input.  We have performed experiments at a finite collection of points $\X = \{x_1, \ldots, x_n\} \subseteq \mathbb{R}^d$. We assume without loss of generality that each element of $\X$ is unique.

We consider two sets of observations:
\begin{itemize}
\item The set of of observations of $y(x)$, $\Y = \{ (x_i, y(x_i)) : i=1,\ldots,n \}$.
\item The set of of observations of $h(x)$, $\H = \{ (x_i, h(x_i)) : i=1,\ldots,n \}$.
\end{itemize}

$h(\cdot)$ is drawn from a Gaussian process with known mean function and kernel.

We then compare two different posterior distributions, one obtained from conditioning on $\Y$ and the other from conditioning on $\H$. Observe that $\Y$ is computable from $\H$ and so conditioning on $\H$ is the same as conditioning on both $\Y$ and $\H$.

Past work has focused on comparing the performance of these two models:

\paragraph{Model 1 (Classical GP regression)} We place a Gaussian process prior on $y(\cdot)$. To perform inference, we calculate the conditional distribution of $y(x)$ at a new location $x$ given $\Y$.

\paragraph{Model 2 (Composite functions)} We place a Gaussian process prior on $h(\cdot)$. This implies a prior distribution on $y(\cdot)$ via the identify $y(x) = g(h(x))$. We calculate conditional distributions given $\H$.

We additionally compare with:

\paragraph{Model 3 (Composite functions prior without composite observations)} As with Model 2, we place a Gaussian process prior on $h(\cdot)$. We calculate conditional distributions given $\Y$.

Both model 3 and model 1 use the same prior. They differ only in the information provided. Thus, we feel that this comparison provides a more meaningful view on the benefits of incorporating composite information into prediction and Gaussian process regresion.

\section{Analysis}

We analyze the difference between model 1 and model 3.  Specifically, we will compare
the cdfs $P(y(x) \le a \mid \Y)$ with $P(y(x) \le a \mid \H)$.

\subsection{Calculating $P(y(x) \le a \mid \Y)$}
We begin by observing that 
$P(y(x) \le a \mid \Y)$
can be computed via Bayes rule.

Given an observation $y=y(x)$ at some $x$, consider the set of values for
$h(x)$ that are compatible with $y$. This set is $H(y) := g^{-1}(y)$.
By selecting one element of $H(y(x_i))$ for each observed $y(x_i)$ in $\Y$,
we can construct a set in the same form of $\H$ that is compatible with $\Y$.
\begin{equation}
\{ (x_i, h_i) : i=1,\ldots,n, h_i \in H(y(x_i)) \}
\end{equation}
We refer to such a set as {\it plausible history}.
There are $\prod_{i=1}^n |H(y(x_i))|$ plausible histories because each plausible history is described by a selection from $H(y(x_i))$ for each $i$.
Let $\Hset$ be the set of all plausible histories.

The density of the posterior distribution at some arbitrary point u can be written, for some arbitrary $x$ and $a$,
\begin{equation*}
P(y(x) \le a \mid \Y) = \sum_{H \in \Hset} P(\H = H) P(y(x) \le a \mid \H = H)
\end{equation*}

\paragraph{Assignment for the SURP students:}
SURP students to update this note to explain how to compute $P(y(x) \le a \mid
\H = H)$ and $P(\H = H)$ {\bf without using AI}.  Assume that you have a
function supplied from elsewhere (BoTorch) that can compute the mean and
variance of $h(x)$ given observations of $h(x_i), i=1,\ldots,n$.

Hint: The derivation doesn't need to be long.  Peter's explanation is 80 words
and less than 500 characters.

In the next meeting, Peter would like someone on the SURP team to walk through
this explanation as if it were a lecture and explain it to everyone, writing
the formulas on the whiteboard, without using any typed notes. Handwritten
notes are ok.

\paragraph{Context:} Peter's goal here is not to understand how to compute
these two quantities. He already knows how to do it. And if he didn't know, he
expects that AI can easily describe the method and write code for it.  The goal
is to help the SURP students understand the notation and ideas above, and how
to do Bayesian inference here, so that they can make progress on the things
that Peter doesn't know how to do, and where AI will not be reliable.




\subsection{Results That We Should Be Able to Show}

It is straightforward to prove that the distributions are equal (i.e., 
$P(y(x) \le a \mid \Y) = P(y(x) \le a \mid \H)$
for all $a$) when $g$ is a bijection.

(This is a good exercise for the SURP students. Do this without AI and give the
proof on the board with only handwritten notes.)

The converse is not true. When $g$ is not a bijection, there are examples where
$P(y(x) \le a \mid \Y) = P(y(x) \le a \mid \H)$.  One trivial example is where
$h(x)$ has a 2-D output, the GP on $h(x)$ is uncorrelated across outputs, $g$
ignores one of the components of $h(x)$, and directly outputs the other one.

Peter believes that when $g$ is linear, the two distributions are equal.
Since many linear functions are not bijections, this is another example showing
a $g$ that is not a bijection where the two distributions are equal.

If $g$ is not a bijection and $y$ is one-dimensional, then in this case it may
be possible to prove that the distributions are not equal.

(These last two paragraphs are also good exercises for the SURP students. They are a little bit harder. But if you can't do it without AI, let's put it off until after the other exercises are done.)

\subsection{Mutual Information Analysis}

Pick some $i$.  Consider the joint probability distribution of $y(x)$ and $h(x_i)$ given that $h(x_i) \in H(y)$ for some fixed $y$. (Later, we will endow $y$ with the distribution of $g(h(x_i))$ under the GP prior.)

Consider the mutual information between $y(x)$ and $h(x_i)$. How large is this mutual information for different choices of $g$?  Let's do a few examples where $y$ is 1-D and some where $y$ is 2-D.

Generally, higher mutual info should correspond to having more value for composite modeling.

\bibliographystyle{plainnat}
\bibliography{refs}

\end{document}
